\begin{review}{Q4.2: AI-suggested five-star implementation self-check}
\draftreview{P3}{P1, P5, P4}{P3 checks corpus preprocessing; P1 checks lexical/dense transformations; P5 checks prompt rendering and output parsing; P4 verifies the clipping analysis.}
\colorbox{yellow!45}{\strut\textbf{DONE FOR DRAFT --- P1}}\par
Approved by Junlin Yin on 6 November 2026 for lexical and Dense preprocessing. Other owners' reviews remain pending.\par\smallskip
\colorbox{yellow!45}{\strut\textbf{DONE FOR DRAFT --- P5}}\par
Approved by Junyu Hu on 6 November 2026 for prompt rendering (eight passages at 1,200 characters) and structured-output parsing. Corpus preprocessing is covered by the P3 box below and the P4 review remains pending; the clipping figures (73.3\%, 33 of 45 development gold passages; 66.9\%, 232 of 347 held-out gold passages) were recomputed by P4 on the final corpus.\par\smallskip
\colorbox{yellow!45}{\strut\textbf{DONE FOR DRAFT --- P3}}\par
P3 scope: stored-text normalization and preserved source placeholders checked. Prompt clipping and parsing remain for P4/P5 verification; the clipping statistic was recomputed by P4 (73.3\%, 33 of 45 development gold passages).\par\smallskip
\aistars{5}\quad\textbf{Provisional AI assessment (7 November 2026, reassessed after preprocessing revision): 5/5 stars.} This is an editorial assessment of Q4.2 requirement coverage, not a course grade or human sign-off. The section explains why structured scientific text needs conservative cleaning rather than extensive microtext normalization. It identifies the transformations applied, the scientific content retained and the distinction between stored passages and retrieval representations. It motivates bounded generation inputs and traceable citation identifiers, explains the critic's shared evidence view, and separates JSON formatting repair from factual validation. The primary context limit and the cross-reference to post-hoc sensitivity analysis preserve the relevant experimental boundary without excessive operational detail. This directly answers whether preprocessing was needed, what was done and why. Context truncation is an explicit design trade-off, not evidence of an incomplete preprocessing answer; an additional content-preservation benchmark is not required by this item. No new experiment, corpus rebuild or owner approval is claimed. Historical owner approvals retain their original scopes.
\end{review}
